\documentclass[11pt,a4paper]{article}

% --- Packages ---
\usepackage[margin=.75in]{geometry}
\usepackage[utf8]{inputenc}
\usepackage[T1]{fontenc}
\usepackage{lmodern} % Fix for font not found error
\usepackage{titlesec}
\usepackage{enumitem}
\usepackage{hyperref}
\usepackage{xcolor}
\usepackage{fancyhdr}
\usepackage{comment}
\usepackage{cleveref}
\usepackage[bottom]{footmisc}
\usepackage{graphicx}
\usepackage{longtable}
\usepackage{booktabs}
\usepackage{array}
\usepackage{calc}
\usepackage{etoolbox}

% Standard Pandoc setup
\providecommand{\tightlist}{%
  \setlength{\itemsep}{0pt}\setlength{\parskip}{0pt}}

\crefformat{section}{\S#2#1#3}
\crefformat{subsection}{\S#2#1#3}
\crefformat{subsubsection}{\S#2#1#3}
\crefformat{figure}{#2Figure~#1#3}
\crefformat{footnote}{#2\footnotemark[#1]#3}

\linespread{0.9} % Riduce lo spazio tra le linee

% --- Colors ---
\definecolor{primary}{RGB}{0, 0, 0}
\definecolor{links}{HTML}{0000EE}
\definecolor{blue}{HTML}{26428B}

\hypersetup{
    colorlinks=true,
    linkcolor=links,
    urlcolor=links,
    citecolor=links,
}

% --- Section Formatting ---
\titleformat{\section}
  {\normalfont\Large\bfseries\color{blue}}{\thesection}{1em}{}
\titleformat{\subsection}
{\normalfont\large\bfseries\color{blue}}{\thesubsection}{1em}{}
\titleformat{\subsubsection}
{\normalfont\bfseries\color{blue}}{\thesubsubsection}{1em}{}

% --- Header and Footer ---
\pagestyle{fancy}
\fancyhead[L]{\small Giuseppe Capasso}
\fancyhead[R]{\small vapt}
\fancyfoot[C]{\thepage}
\renewcommand{\headrulewidth}{0.4pt}

% --- Custom Title Command ---
\newcommand{\proposaltitle}[1]{
    \begin{center}
        \vspace*{0.1cm}
        {\LARGE \bfseries \color{blue} #1} \\[1.5ex]
        {\large \textbf{Giuseppe Capasso}} \\[1ex]
                \small{
            \vspace{0.1cm}
            \textbf{Materia:} vapt\\
        }
            \end{center}
    \vspace{0.3cm}
    \hrule height 0.5pt
    \vspace{0.5cm}
}

\begin{document}

\proposaltitle{Dispensa Lezione 2: Web Fuzzing, Shell e Attacchi in Laboratorio}

\section{Dispensa Lezione 2: Web Fuzzing, Shell e Attacchi in
Laboratorio}\label{dispensa-lezione-2-web-fuzzing-shell-e-attacchi-in-laboratorio}

\subsection{1. Introduzione alla
Lezione}\label{introduzione-alla-lezione}

Questa seconda lezione è divisa in due fasi. Nella prima parte
analizzeremo insieme la compromissione completa di due macchine target
sulla piattaforma HackTheBox. Vedremo come combinare
l\textquotesingle enumerazione dei servizi web, lo sfruttamento di
vulnerabilità specifiche e l\textquotesingle ottenimento di privilegi di
amministratore.

Nella seconda parte, vi dividerete in due Red Team autonomi per
attaccare due ambienti distinti: il primo focalizzato sulla sicurezza
Web (DVWA) e il secondo sulla sicurezza di sistema (Metasploitable). Al
termine, ogni team dovrà rendicontare il proprio lavoro
all\textquotesingle altro gruppo.

\subsection{2. Strumenti e Concetti (Fase
Dimostrativa)}\label{strumenti-e-concetti-fase-dimostrativa}

Durante la dimostrazione pratica, utilizzeremo i seguenti strumenti e
concetti fondamentali per il Penetration Testing moderno.

\subsubsection{2.1 Enumeration: Nmap e
FFUF}\label{enumeration-nmap-e-ffuf}

La fase di ricognizione è il cuore dell\textquotesingle attacco.

Nmap: Lo utilizzeremo per mappare le porte aperte e identificare la
superficie di attacco.

FFUF (Fuzz Faster U Fool): Quando identifichiamo un server web (es.
porta 80 o 443), non basta navigare il sito tramite browser. FFUF è un
web fuzzer scritto in Go, estremamente veloce, che ci permette di
scoprire directory nascoste, file non linkati e parametri invisibili
testando migliaia di parole (prese da una wordlist) contro
l\textquotesingle URL target.

Sintassi base: ffuf -w /usr/share/wordlists/dirb/common.txt -u
http://\textless IP\_TARGET\textgreater/FUZZ

La parola FUZZ nell\textquotesingle URL agisce da segnaposto: FFUF la
sostituirà con ogni parola della wordlist.

\subsubsection{2.2 Exploitation: Il protocollo
MCP}\label{exploitation-il-protocollo-mcp}

Nel nostro scenario d\textquotesingle esempio,
l\textquotesingle enumerazione rivelerà l\textquotesingle utilizzo del
protocollo MCP. Esamineremo come l\textquotesingle implementazione o la
configurazione di determinati protocolli (spesso legati ad applicazioni
di terze parti o sistemi di gestione) possa presentare vulnerabilità
sfruttabili per l\textquotesingle esecuzione di codice remoto (RCE).
L\textquotesingle obiettivo è capire che qualsiasi servizio esposto, non
solo i classici web o database, rappresenta un potenziale vettore di
ingresso.

2.3 Post-Exploitation: Bind Shell vs Reverse Shell

Una volta trovata una vulnerabilità che permette
l\textquotesingle esecuzione di comandi (RCE), il nostro obiettivo è
ottenere una Shell interattiva sul sistema bersaglio. Esistono due
tipologie principali:

Bind Shell: Il payload malevolo apre una porta in ascolto sulla macchina
bersaglio (es. porta 4444) e lega a essa una shell (come /bin/bash o
cmd.exe). L\textquotesingle attaccante si connette poi dal proprio
computer a quella porta.

Svantaggio: I firewall moderni bloccano quasi sempre le connessioni in
ingresso su porte non standard.

Reverse Shell: È la tecnica più utilizzata. L\textquotesingle attaccante
apre una porta in ascolto sulla propria macchina (tramite Netcat: nc
-lvnp 9001). Il payload eseguito sulla macchina bersaglio stabilisce una
connessione in uscita verso il computer dell\textquotesingle attaccante,
consegnandogli la shell.

Vantaggio: I firewall sono generalmente molto più permissivi sul
traffico in uscita, permettendo al payload di "chiamare casa" senza
essere bloccato.

\subsubsection{2.4 Privilege Escalation:
LinPEAS}\label{privilege-escalation-linpeas}

Ottenuto l\textquotesingle accesso iniziale, spesso ci ritroveremo con
privilegi bassi (es. utente www-data). Per diventare amministratori
(root), dobbiamo cercare falle interne al sistema.

LinPEAS (Linux Privilege Escalation Awesome Script) è uno script bash
che automatizza la ricerca di vettori di Privilege Escalation. Analizza
permessi errati, file SUID, cronjob mal configurati, password in chiaro
nei log e kernel vulnerabili, evidenziando i risultati con un sistema a
colori per facilitarne la lettura.

\subsection{3. Attività di Laboratorio (Red
Teaming)}\label{attivituxe0-di-laboratorio-red-teaming}

La classe viene ora divisa in due gruppi. Ogni gruppo ha
l\textquotesingle obiettivo di compromettere il proprio target, fare più
danni possibili (simulando un attacco reale), documentare i passaggi e
preparare una presentazione tecnica per i colleghi.

\paragraph{Gruppo 1: Web Application Hacking
(DVWA)}\label{gruppo-1-web-application-hacking-dvwa}

Target: Damn Vulnerable Web App (DVWA)

Focus: OWASP Top 10, SQL Injection, Command Injection, XSS, File Upload.

Setup (Docker Compose): Creerete un ambiente effimero tramite container.

Create una cartella: mkdir lab\_dvwa \&\& cd lab\_dvwa

Create un file docker-compose.yml e inserite le istruzioni standard per
l\textquotesingle immagine di DVWA (es. vulnerables/web-dvwa).

Avviate l\textquotesingle infrastruttura con il comando:

Bash

docker compose up -d

Accedete all\textquotesingle IP locale di Kali sulla porta 80 tramite
browser. Login di default: admin / password.

Obiettivo: Affrontare le sfide presenti nel menu laterale, partendo dal
livello di sicurezza "Low". Cercate di estrarre il database tramite SQLi
e di ottenere una reverse shell tramite Command Injection. Consultate
guide online e la documentazione OWASP in caso di blocco.

Ripristino Ambiente: Se compromettete il sistema al punto da renderlo
instabile, potete resettarlo azzerando il container:

Bash

docker compose down

docker compose up -d

\paragraph{Gruppo 2: Infrastructure \& System Hacking
(Metasploitable)}\label{gruppo-2-infrastructure-system-hacking-metasploitable}

Target: Metasploitable 2

Focus: Enumerazione dei servizi, sfruttamento di software obsoleto (FTP,
Samba, distccd, SSH), Privilege Escalation.

Setup (Snapshot di Sicurezza): La macchina è già configurata su
VirtualBox. Poiché il vostro obiettivo è compromettere a fondo il
sistema (modificando file di configurazione, aggiungendo utenti root,
ecc.), è obbligatorio creare uno Snapshot prima di iniziare.

Su VirtualBox, selezionate la VM Metasploitable -\textgreater{}
"Istantanee" -\textgreater{} "Cattura".

Questo vi permetterà di riportare la macchina allo stato originario con
un click alla fine dell\textquotesingle esercitazione o nel caso in cui
"rompiate" il sistema operativo.

Obiettivo: Partite da zero con Nmap. Identificate i servizi in ascolto e
cercate exploit noti (potete usare searchsploit su Kali o cercare le CVE
online). Puntate a ottenere l\textquotesingle accesso come utente root
in almeno 3 modi diversi sfruttando 3 servizi differenti.

\end{document}
